\chapter{Hilbert空间的几何}\label{chap:I-04}
\FAChapterOpening
{建立复内积第一变量线性约定下的 Hilbert 几何；证明闭凸投影、闭子空间正交分解、Hilbert 空间 Riesz 表示、Bessel 不等式、Parseval 恒等式与正交基完备性判据；在 \(\displaystyle \ell^2\) 与 \(\displaystyle L^2\) 中核验标准 Hilbert 结构.}
{I-03，特别是\autoref{fa:BAN-A01}与\autoref{i03:Lp-banach}.}
{\begin{center}
\begin{tikzpicture}[x=2.35cm,y=1.15cm]
\node[fa/node] (ban) at (0,0) {赋范/Banach空间};
\node[fa/node] (c01) at (1,0) {正交计算};
\node[fa/node] (a01) at (2,0) {Hilbert投影定理};
\node[fa/node] (a02) at (3,0) {Riesz表示定理};
\node[fa/node] (b01) at (1.7,-1) {Bessel不等式/Parseval恒等式};
\node[fa/node] (b02) at (2.8,-1) {正交基判据};
\draw[fa/arrow] (ban) -- (c01);
\draw[fa/arrow] (c01) -- (a01);
\draw[fa/arrow] (a01) -- (a02);
\draw[fa/arrow] (c01) -- (b01);
\draw[fa/arrow] (a01) -- (b01);
\draw[fa/arrow] (b01) -- (b02);
\end{tikzpicture}
\end{center}}

\section{Hilbert与正交性}\label{sec:i04-hilbert-orthogonality}
\index{inner product@内积}\index{Hilbert space@Hilbert空间}\index{orthogonality@正交}\index{orthogonal complement@正交补}

\subsection{内积与诱导范数}

全文固定 \(\displaystyle \mathbb K\in\{\mathbb R,\mathbb C\}\). 复数域上的内积始终第一变量线性、第二变量共轭线性.

\begin{FADefinition}[B][内积与诱导长度]{i04:inner-product}
设 \(\displaystyle H\) 为 \(\displaystyle \mathbb K\)-线性空间. 映射 \(\displaystyle \langle\cdot,\cdot\rangle:H\times H\to\mathbb K\) 称为内积，若对任意 \(\displaystyle x,x_1,x_2,y,y_1,y_2\in H\) 与 \(\displaystyle \alpha,\beta\in\mathbb K\) 有 \(\displaystyle \langle\alpha x_1+\beta x_2,y\rangle =\alpha\langle x_1,y\rangle+\beta\langle x_2,y\rangle\)，
\[
\langle x,\alpha y_1+\beta y_2\rangle
=\overline\alpha\langle x,y_1\rangle+\overline\beta\langle x,y_2\rangle,
\]
以及 \(\displaystyle \langle x,y\rangle=\overline{\langle y,x\rangle}\)、\(\displaystyle \langle x,x\rangle\geqslant0\) 且 \(\displaystyle \langle x,x\rangle=0\iff x=0\). 定义诱导长度 \(\displaystyle \|x\|:=\sqrt{\langle x,x\rangle}\).
当 \(\displaystyle \mathbb K=\mathbb R\) 时称为实内积空间；当 \(\displaystyle \mathbb K=\mathbb C\) 时称为复内积空间. 实数域情形中复共轭恒等于自身.
\end{FADefinition}

\begin{FAProposition}[C][Hilbert正交计算工具]{fa:HIL-C01}
在\autoref{i04:inner-product}的条件下成立：
\begin{enumerate}[label=\textnormal{(\roman*)}]
\item Cauchy--Schwarz 不等式 \(\displaystyle |\langle x,y\rangle|\leqslant\|x\|\|y\|\)，且等号成立当且仅当 \(\displaystyle x,y\) 线性相关；
\item \(\displaystyle \|\cdot\|\) 满足\autoref{fa:BAN-A01}的全部范数公理；
\item 若 \(\displaystyle \langle x,y\rangle=0\)，则 \(\displaystyle \|x+y\|^2=\|x\|^2+\|y\|^2\)；
\item \(\displaystyle \|x+y\|^2+\|x-y\|^2=2\|x\|^2+2\|y\|^2\)；
\item 若 \(\displaystyle x_1,\dots,x_n\) 两两正交，则 \(\displaystyle \|\sum_{k=1}^n x_k\|^2=\sum_{k=1}^n\|x_k\|^2\)；
\item 对任意线性子空间 \(\displaystyle M\leqslant H\)，其正交补 \(\displaystyle M^\perp\) 是闭线性子空间.
\end{enumerate}
此外，实情形与复情形分别有极化恒等式
\[
\langle x,y\rangle
=\frac14\left(\|x+y\|^2-\|x-y\|^2\right),
\quad \mathbb K=\mathbb R,
\]
与
\[
\langle x,y\rangle
=\frac14\left(\|x+y\|^2-\|x-y\|^2\right)
+\frac{\mathrm{i}}4\left(\|x+\mathrm{i}y\|^2-\|x-\mathrm{i}y\|^2\right),
\quad \mathbb K=\mathbb C.
\]
\end{FAProposition}

\begin{FAProof}
先证 Cauchy--Schwarz. 若 \(\displaystyle y=0\)，则 \(\displaystyle \langle x,y\rangle=0\)，不等式成立. 设 \(\displaystyle y\neq0\)，取 \(\displaystyle \lambda=\frac{\langle x,y\rangle}{\|y\|^2}\).
由第一变量线性与第二变量共轭线性，
\begin{align*}
\displaystyle 0&\leqslant\|x-\lambda y\|^2=\langle x-\lambda y,x-\lambda y\rangle\\
&\displaystyle =\|x\|^2-\langle x,\lambda y\rangle-\langle\lambda y,x\rangle+\langle\lambda y,\lambda y\rangle\\
&\displaystyle =\|x\|^2-\overline\lambda\langle x,y\rangle-\lambda\overline{\langle x,y\rangle}+|\lambda|^2\|y\|^2\\[6pt]
&\displaystyle =\|x\|^2-\frac{|\langle x,y\rangle|^2}{\|y\|^2}.
\end{align*}
故 \(\displaystyle |\langle x,y\rangle|^2\leqslant\|x\|^2\|y\|^2\)，取非负平方根即得不等式.

若 \(\displaystyle x=0\) 或 \(\displaystyle y=0\)，两向量线性相关且等号成立. 现设 \(\displaystyle x,y\neq0\). 上述恒等式表明 Cauchy--Schwarz 取等号当且仅当 \(\displaystyle \|x-\lambda y\|^2=0\)，即 \(\displaystyle x=\lambda y\). 反之，若 \(\displaystyle x=\alpha y\)，则 \(\displaystyle |\langle x,y\rangle|=|\alpha|\|y\|^2=\|x\|\|y\|\). 因而等号条件闭合.

验证范数公理. 非负性与正定性直接来自内积正定性. 对 \(\displaystyle \alpha\in\mathbb K\)，
\[
\|\alpha x\|^2
=\langle\alpha x,\alpha x\rangle
=\alpha\overline\alpha\langle x,x\rangle
=|\alpha|^2\|x\|^2,
\]
故 \(\displaystyle \|\alpha x\|=|\alpha|\|x\|\). 再由 Cauchy--Schwarz，
\begin{align*}
\|x+y\|^2
&=\|x\|^2+\langle x,y\rangle+\langle y,x\rangle+\|y\|^2=\|x\|^2+2\operatorname{Re}\langle x,y\rangle+\|y\|^2\\
&\leqslant\|x\|^2+2|\langle x,y\rangle|+\|y\|^2\\
&\leqslant(\|x\|+\|y\|)^2.
\end{align*}
两端取非负平方根得到三角不等式，故 \(\displaystyle \|\cdot\|\) 确为范数.

若 \(\displaystyle \langle x,y\rangle=0\)，则共轭对称性给出 \(\displaystyle \langle y,x\rangle=0\)，因此 \(\displaystyle \|x+y\|^2 =\langle x+y,x+y\rangle =\|x\|^2+\|y\|^2\).
同样逐项展开得
\[
\|x+y\|^2+\|x-y\|^2=2\langle x,x\rangle+2\langle y,y\rangle=2\|x\|^2+2\|y\|^2.
\]
若 \(\displaystyle x_1,\dots,x_n\) 两两正交，则
\[
\left\|\sum_{k=1}^{n}x_k\right\|^2
=\sum_{j=1}^{n}\sum_{k=1}^{n}\langle x_j,x_k\rangle
=\sum_{k=1}^{n}\|x_k\|^2,
\]
因为 \(\displaystyle j\neq k\Rightarrow\langle x_j,x_k\rangle=0\).

设 \(\displaystyle z_1,z_2\in M^\perp\)、\(\displaystyle \alpha,\beta\in\mathbb K\). 对每个 \(\displaystyle m\in M\)，
\[
\langle\alpha z_1+\beta z_2,m\rangle
=\alpha\langle z_1,m\rangle+\beta\langle z_2,m\rangle=0,
\]
故 \(\displaystyle M^\perp\) 为线性子空间. 若 \(\displaystyle z_n\in M^\perp\) 且 \(\displaystyle z_n\to z\)，则对每个 \(\displaystyle m\in M\)，Cauchy--Schwarz 给出 \(\displaystyle |\langle z,m\rangle| =|\langle z-z_n,m\rangle| \leqslant\|z-z_n\|\|m\|\to0\).
因此 \(\displaystyle \langle z,m\rangle=0\)，故 \(\displaystyle z\in M^\perp\)，从而 \(\displaystyle M^\perp\) 闭.

最后核验极化恒等式. 由共轭对称性，
\[
\|x+y\|^2-\|x-y\|^2
=2\langle x,y\rangle+2\langle y,x\rangle
=4\operatorname{Re}\langle x,y\rangle.
\]
实情形中右端即 \(\displaystyle4\langle x,y\rangle\). 复情形中，第二变量共轭线性给出
\begin{align*}
\|x+\mathrm{i}y\|^2
&=\|x\|^2+\|y\|^2-\mathrm{i}\langle x,y\rangle+\mathrm{i}\overline{\langle x,y\rangle},\\
\|x-\mathrm{i}y\|^2
&=\|x\|^2+\|y\|^2+\mathrm{i}\langle x,y\rangle-\mathrm{i}\overline{\langle x,y\rangle}.
\end{align*}
两式相减得到 \(\displaystyle \|x+\mathrm{i}y\|^2-\|x-\mathrm{i}y\|^2 =4\operatorname{Im}\langle x,y\rangle\).
与实部公式合并即得到所列复极化恒等式.\hfill\(\square\)

\end{FAProof}

\begin{FADefinition}[B][Hilbert空间、正交族与正交补]{i04:hilbert-orthogonal-definitions}
内积空间 \(\displaystyle H\) 若对诱导范数完备，则称为 Hilbert 空间. 对 \(\displaystyle x,y\in H\)，写 \(\displaystyle x\perp y\iff\langle x,y\rangle=0\). 族 \(\displaystyle (x_j)_{j\in J}\) 称为正交族，若 \(\displaystyle j\neq k\Rightarrow x_j\perp x_k\)；若另有 \(\displaystyle \|x_j\|=1\)，则称为标准正交族. 对 \(\displaystyle A\subseteq H\)，定义 \(\displaystyle A^\perp:=\{x\in H:\forall\,a\in A,\ \langle x,a\rangle=0\}\).
对 \(\displaystyle E\subseteq H\)，记 \(\displaystyle \overline{\operatorname{span}}E\) 为其线性张成的闭包.
\end{FADefinition}

\newpage
\subsection{\texorpdfstring{\(\displaystyle\ell^2\)}{ℓ^2}的Hilbert结构}

\begin{FAProposition}[B][\(\displaystyle\ell^2\)的标准内积]{i04:ell2-hilbert}
对 \(\displaystyle x=(x_n),y=(y_n)\in\ell^2\)，定义 \(\displaystyle \langle x,y\rangle_{\ell^2} :=\sum_{n=1}^{\infty}x_n\overline{y_n}\).
该级数绝对收敛并给出 \(\displaystyle \ell^2\) 上的内积；其诱导范数正是 I-03 中的 \(\displaystyle \|x\|_2=(\sum_{n=1}^{\infty}|x_n|^2)^{\frac{1}{2}}\). 因而 \(\displaystyle \ell^2\) 是 Hilbert 空间.
\end{FAProposition}

\begin{FAProof}
对每个 \(\displaystyle n\)，由 \(\displaystyle 2|x_n||y_n|\leqslant|x_n|^2+|y_n|^2\)，
\[
\sum_{n=1}^{\infty}|x_n\overline{y_n}|
\leqslant\frac12\sum_{n=1}^{\infty}|x_n|^2
+\frac12\sum_{n=1}^{\infty}|y_n|^2
<\infty.
\]
故定义有意义. 绝对收敛允许逐项进行有限线性组合，因而第一变量线性与第二变量共轭线性成立；共轭对称性来自
\[
\overline{\langle y,x\rangle_{\ell^2}}
=\overline{\sum_{n=1}^{\infty}y_n\overline{x_n}}
=\sum_{n=1}^{\infty}x_n\overline{y_n}
=\langle x,y\rangle_{\ell^2}.
\]
又 \(\displaystyle \langle x,x\rangle_{\ell^2}=\sum_{n=1}^{\infty}|x_n|^2\geqslant0\)，且该和为零当且仅当每个 \(\displaystyle x_n=0\). 因而这确为内积，且诱导范数与 I-03 的 \(\displaystyle \ell^2\) 范数相同. 由\autoref{i03:Lp-banach}在计数测度情形的结论，\(\displaystyle \ell^2\) 对此范数完备，所以是 Hilbert 空间.\hfill\(\square\)

\end{FAProof}

\begin{FAExample}[序列Hilbert空间]
\autoref{i04:ell2-hilbert}将 I-03 的 Banach 空间 \(\displaystyle \ell^2\) 升格为 Hilbert 空间，但未改变其集合或范数.
\end{FAExample}

\subsection{\texorpdfstring{\(\displaystyle L^2(\Omega)\)}{L^2}的Hilbert结构}

\begin{FAProposition}[B][\(\displaystyle L^2(\Omega)\)的标准内积]{i04:L2-hilbert}
设 \(\displaystyle (\Omega,\Sigma,\mu)\) 为测度空间. 对 \(\displaystyle f,g\in L^2(\Omega)\) 定义 \(\displaystyle \langle f,g\rangle_{L^2} :=\int_\Omega f\overline g\,\mathrm{d}\mu\).
该表达式对几乎处处等价类良定义并给出 \(\displaystyle L^2(\Omega)\) 上的内积；诱导范数恰为 I-03 的 \(\displaystyle L^2\) 范数. 因而 \(\displaystyle L^2(\Omega)\) 是 Hilbert 空间.
\end{FAProposition}

\begin{FAProof}
取任意可测代表元，逐点有 \(\displaystyle 2|f\overline g|\leqslant|f|^2+|g|^2\). 因 \(\displaystyle |f|^2,|g|^2\in L^1(\Omega)\)，得到 \(\displaystyle f\overline g\in L^1(\Omega)\)，故积分有限. 若 \(\displaystyle f=f'\) 几乎处处且 \(\displaystyle g=g'\) 几乎处处，则在两个零测例外集之并的补集上 \(\displaystyle f\overline g=f'\overline{g'}\)，故积分相同；因此内积与代表元选择无关.

积分线性与复共轭的基本性质给出第一变量线性、第二变量共轭线性及共轭对称性. 又 \(\displaystyle \langle f,f\rangle_{L^2} =\int_\Omega|f|^2\,\mathrm{d}\mu \geqslant0\)，
且积分为零当且仅当 \(\displaystyle f=0\) 几乎处处，即 \(\displaystyle f\) 是零等价类. 因而这确为内积，且
\[
\sqrt{\langle f,f\rangle_{L^2}}
=\left(\int_\Omega|f|^2\,\mathrm{d}\mu\right)^{\frac{1}{2}}
=\|f\|_2.
\]
由\autoref{i03:Lp-banach}在 \(\displaystyle p=2\) 时的完备性，\(\displaystyle L^2(\Omega)\) 是 Hilbert 空间.\hfill\(\square\)

\end{FAProof}

\begin{FAExample}[函数Hilbert空间]
本章只使用一般 \(\displaystyle L^2(\Omega)\) 的 Hilbert 结构，不宣称任何未经完备性证明的具体函数系为 \(\displaystyle L^2\) 标准正交基.
\end{FAExample}

\section{最佳逼近}\label{sec:i04-best-approximation}
\index{best approximation@最佳逼近}\index{projection theorem@投影定理}

\subsection{距离与闭凸投影}

对非空集合 \(\displaystyle C\subseteq H\) 与 \(\displaystyle x\in H\)，定义 \(\displaystyle d(x,C):=\inf_{z\in C}\|x-z\|\).
若存在唯一 \(\displaystyle y\in C\) 满足 \(\displaystyle \|x-y\|=d(x,C)\)，记 \(\displaystyle y=\mathcal P_Cx\).

\begin{FATheorem}[A][Hilbert投影定理]{fa:HIL-A01}
设 \(\displaystyle H\) 为 Hilbert 空间，\(\displaystyle C\subseteq H\) 非空、闭、凸. 则 \(\displaystyle \forall\,x\in H\ \exists\,!y\in C:\ \|x-y\|=d(x,C)\).
亦即 \(\displaystyle \mathcal P_Cx\) 对每个 \(\displaystyle x\in H\) 唯一存在.
\end{FATheorem}

\begin{FAProof}
固定 \(\displaystyle x\in H\)，令 \(\displaystyle d=d(x,C)\). 由下确界定义，对每个 \(\displaystyle n\in\mathbb N\) 可取 \(\displaystyle y_n\in C\) 使 \(\displaystyle d\leqslant\|x-y_n\|<\sqrt{d^2+\frac{1}{n}}\)，故 \(\displaystyle \|x-y_n\|^2\to d^2\).

对 \(\displaystyle n,m\in\mathbb N\)，凸性给出 \(\displaystyle \frac{y_n+y_m}{2}\in C\)，从而 \(\displaystyle \left\|x-\frac{y_n+y_m}{2}\right\|\geqslant d\).
对 \(\displaystyle a=x-y_n\)、\(\displaystyle b=x-y_m\) 使用平行四边形恒等式，得到
\begin{align*}
\displaystyle \|y_n-y_m\|^2
&\displaystyle =\|a-b\|^2=2\|a\|^2+2\|b\|^2-\|a+b\|^2\\[6pt]
&\displaystyle =2\|x-y_n\|^2+2\|x-y_m\|^2-4\left\|x-\frac{y_n+y_m}{2}\right\|^2\\[6pt]
&\displaystyle \leqslant2\|x-y_n\|^2+2\|x-y_m\|^2-4d^2.
\end{align*}
当 \(\displaystyle n,m\to\infty\) 时右端趋于零，故 \(\displaystyle (y_n)\) 为 Cauchy 列. Hilbert 完备性给出某个 \(\displaystyle y\in H\) 使 \(\displaystyle y_n\to y\). 因 \(\displaystyle C\) 闭，\(\displaystyle y\in C\). 范数的反三角不等式给出 \(\displaystyle \bigl|\|x-y_n\|-\|x-y\|\bigr| \leqslant\|y_n-y\|\to0\)，
故 \(\displaystyle \|x-y\|=d\)，存在性成立.

若 \(\displaystyle y_1,y_2\in C\) 且 \(\displaystyle \|x-y_1\|=\|x-y_2\|=d\)，则 \(\displaystyle \frac{y_1+y_2}{2}\in C\). 再用同一平行四边形恒等式，
\[
\|y_1-y_2\|^2=2d^2+2d^2-4\left\|x-\frac{y_1+y_2}{2}\right\|^2\leqslant4d^2-4d^2=0.
\]
故 \(\displaystyle y_1=y_2\)，唯一性成立.\hfill\(\square\)

\end{FAProof}

\subsection{投影的变分刻画}

\begin{FAProposition}[B][闭凸投影的变分刻画]{i04:projection-variational}
在\autoref{fa:HIL-A01}的条件下，对 \(\displaystyle x\in H\)、\(\displaystyle y\in C\)，有 \(\displaystyle y=\mathcal P_Cx \iff \forall\,z\in C,\ \operatorname{Re}\langle x-y,z-y\rangle\leqslant0\).
\end{FAProposition}

\begin{FAProof}
设 \(\displaystyle y=\mathcal P_Cx\). 固定 \(\displaystyle z\in C\). 对 \(\displaystyle t\in[0,1]\)，凸性给出 \(\displaystyle y+t(z-y)\in C\)，于是 \(\displaystyle \|x-y\|^2 \leqslant\|x-y-t(z-y)\|^2\).
令 \(\displaystyle a=x-y\)、\(\displaystyle b=z-y\). 因 \(\displaystyle t\in\mathbb R\)，逐项展开得到 \(\displaystyle \|a-tb\|^2 =\|a\|^2-2t\operatorname{Re}\langle a,b\rangle+t^2\|b\|^2\).
故对每个 \(\displaystyle t\in(0,1]\)， \(\displaystyle 2\operatorname{Re}\langle x-y,z-y\rangle \leqslant t\|z-y\|^2\).
令 \(\displaystyle t\downarrow0\)，得到 \(\displaystyle \operatorname{Re}\langle x-y,z-y\rangle\leqslant0\).

反之，设该实部不等式对全部 \(\displaystyle z\in C\) 成立. 则
\begin{align*}
\|x-z\|^2
&=\|x-y-(z-y)\|^2=\|x-y\|^2+\|z-y\|^2-2\operatorname{Re}\langle x-y,z-y\rangle\\
&\geqslant\|x-y\|^2.
\end{align*}
故 \(\displaystyle y\) 达到 \(\displaystyle d(x,C)\). 由\autoref{fa:HIL-A01}的唯一性，\(\displaystyle y=\mathcal P_Cx\).\hfill\(\square\)

\end{FAProof}

\newpage
\section{投影与正交分解}\label{sec:i04-orthogonal-decomposition}
\index{orthogonal projection@正交投影}\index{orthogonal decomposition@正交分解}\index{double orthogonal complement@双正交补}

\subsection{闭子空间投影}

\begin{FAProposition}[B][闭子空间投影与正交分解]{i04:closed-subspace-projection}
设 \(\displaystyle M\leqslant H\) 为闭线性子空间. 对任意 \(\displaystyle x\in H\)、\(\displaystyle y\in M\)， \(\displaystyle y=\mathcal P_Mx \iff y\in M\ \text{且}\ x-y\in M^\perp\).
并且 \(\displaystyle H=M\oplus M^\perp\).
定义 \(\displaystyle \mathcal P_Mx\) 为分解 \(\displaystyle x=m+n\)、\(\displaystyle m\in M\)、\(\displaystyle n\in M^\perp\) 中的 \(\displaystyle m\). 则 \(\displaystyle \mathcal P_M\) 线性，且 \(\displaystyle \|\mathcal P_Mx\|\leqslant\|x\|\).
\end{FAProposition}

\begin{FAProof}
闭线性子空间 \(\displaystyle M\) 非空且凸，故\autoref{fa:HIL-A01}适用. 若 \(\displaystyle y=\mathcal P_Mx\)，由\autoref{i04:projection-variational}，对每个 \(\displaystyle w\in M\) 可分别取 \(\displaystyle z=y+w\) 与 \(\displaystyle z=y-w\)，得到
\[
\operatorname{Re}\langle x-y,w\rangle\leqslant0,
\quad
-\operatorname{Re}\langle x-y,w\rangle\leqslant0.
\]
故 \(\displaystyle \operatorname{Re}\langle x-y,w\rangle=0\). 若 \(\displaystyle \mathbb K=\mathbb C\)，再以 \(\displaystyle \mathrm{i}w\) 与 \(\displaystyle -\mathrm{i}w\) 代入，得到 \(\displaystyle \operatorname{Re}\langle x-y,\mathrm{i}w\rangle=0\). 第二变量共轭线性给出 \(\displaystyle \langle x-y,\mathrm{i}w\rangle =-\mathrm{i}\langle x-y,w\rangle\)，
故 \(\displaystyle \operatorname{Im}\langle x-y,w\rangle=0\). 实、复两种情形均得到 \(\displaystyle \langle x-y,w\rangle=0\)，即 \(\displaystyle x-y\in M^\perp\).

反之，若 \(\displaystyle y\in M\) 且 \(\displaystyle x-y\in M^\perp\)，则对每个 \(\displaystyle z\in M\)，\(\displaystyle z-y\in M\)，故 \(\displaystyle \operatorname{Re}\langle x-y,z-y\rangle=0\leqslant0\). 由\autoref{i04:projection-variational}，\(\displaystyle y=\mathcal P_Mx\).

对任意 \(\displaystyle x\in H\)，令 \(\displaystyle m=\mathcal P_Mx\)、\(\displaystyle n=x-m\). 上述刻画给出 \(\displaystyle m\in M\)、\(\displaystyle n\in M^\perp\)，故 \(\displaystyle x=m+n\). 若另有 \(\displaystyle x=m'+n'\)，则 \(\displaystyle m-m'=n'-n\in M\cap M^\perp\). 对 \(\displaystyle v\in M\cap M^\perp\)，\(\displaystyle \|v\|^2=\langle v,v\rangle=0\)，故 \(\displaystyle v=0\). 因而分解唯一，\(\displaystyle H=M\oplus M^\perp\).

若 \(\displaystyle x=m+n\)、\(\displaystyle x'=m'+n'\) 为对应正交分解，则对 \(\displaystyle \alpha,\beta\in\mathbb K\)， \(\displaystyle \alpha x+\beta x'=(\alpha m+\beta m')+(\alpha n+\beta n')\),
且前一项属于 \(\displaystyle M\)，后一项属于 \(\displaystyle M^\perp\). 由分解唯一性，\(\displaystyle \mathcal P_M(\alpha x+\beta x')=\alpha\mathcal P_Mx+\beta\mathcal P_Mx'\)，故 \(\displaystyle \mathcal P_M\) 线性. 最后由勾股恒等式， \(\displaystyle \|x\|^2 =\|\mathcal P_Mx\|^2+\|x-\mathcal P_Mx\|^2 \geqslant\|\mathcal P_Mx\|^2\),
故 \(\displaystyle \|\mathcal P_Mx\|\leqslant\|x\|\).\hfill\(\square\)

\end{FAProof}

\subsection{双正交补}

\begin{FAProposition}[B][双正交补定理]{i04:double-orthogonal-complement}
对任意线性子空间 \(\displaystyle M\leqslant H\)，不要求 \(\displaystyle M\) 闭，有 \(\displaystyle M^{\perp\perp}=\overline M\).
\end{FAProposition}

\begin{FAProof}
先证 \(\displaystyle \overline M\subseteq M^{\perp\perp}\). 取 \(\displaystyle x\in\overline M\) 与 \(\displaystyle z\in M^\perp\). 存在 \(\displaystyle m_n\in M\) 使 \(\displaystyle m_n\to x\). 因 \(\displaystyle \langle m_n,z\rangle=0\)，Cauchy--Schwarz 给出 \(\displaystyle |\langle x,z\rangle| =|\langle x-m_n,z\rangle| \leqslant\|x-m_n\|\|z\|\to0\).
故 \(\displaystyle \langle x,z\rangle=0\). 因 \(\displaystyle z\in M^\perp\) 任意，\(\displaystyle x\in M^{\perp\perp}\).

反向取 \(\displaystyle x\notin\overline M\). 子空间 \(\displaystyle \overline M\) 闭，故由\autoref{i04:closed-subspace-projection}唯一分解 \(\displaystyle x=m+n, \; m\in\overline M, \; n\in(\overline M)^\perp\).
若 \(\displaystyle n=0\)，则 \(\displaystyle x=m\in\overline M\)，矛盾，故 \(\displaystyle n\neq0\). 因 \(\displaystyle M\subseteq\overline M\)，有 \(\displaystyle n\in M^\perp\). 又 \(\displaystyle \langle x,n\rangle =\langle m,n\rangle+\langle n,n\rangle =\|n\|^2>0\)，
所以 \(\displaystyle x\notin M^{\perp\perp}\). 因而 \(\displaystyle M^{\perp\perp}\subseteq\overline M\)，两边相等.\hfill\(\square\)

\end{FAProof}

\subsection{有限正交投影与最小二乘}

\begin{FAApplication}[有限维正交最小二乘]
设 \(\displaystyle e_1,\dots,e_n\) 标准正交，\(\displaystyle M=\operatorname{span}\{e_1,\dots,e_n\}\). 对 \(\displaystyle x\in H\)，令 \(\displaystyle m_0=\sum_{k=1}^{n}\langle x,e_k\rangle e_k\).
对每个 \(\displaystyle j\in\{1,\dots,n\}\)，
\[
\langle x-m_0,e_j\rangle
=\langle x,e_j\rangle
-\sum_{k=1}^{n}\langle x,e_k\rangle\langle e_k,e_j\rangle
=0.
\]
故 \(\displaystyle x-m_0\in M^\perp\)，由\autoref{i04:closed-subspace-projection}， \(\displaystyle \mathcal P_Mx =\sum_{k=1}^{n}\langle x,e_k\rangle e_k\).
对任意 \(\displaystyle m=\sum_{k=1}^{n}c_ke_k\in M\)，正交分解给出 \(\displaystyle \|x-m\|^2 =\|x-\mathcal P_Mx\|^2 +\sum_{k=1}^{n}|c_k-\langle x,e_k\rangle|^2\).
因此 \(\displaystyle m=\mathcal P_Mx\) 是 \(\displaystyle \|x-m\|\)、\(\displaystyle m\in M\) 的唯一最小化元.
\end{FAApplication}

\begin{FARemark}[无限维闭单位球非紧性沿用I-03]
在 \(\displaystyle \ell^2\) 中，标准单位向量满足 \(\displaystyle n\neq m\Rightarrow\|e_n-e_m\|_2=\sqrt2\). 因而 Hilbert 空间的闭单位球仍可非紧；这正是 I-03 中\autoref{fa:BAN-A02}所建立的无限维闭有界集非紧机制在 \(\displaystyle \ell^2\) 中的具体表现，本章不创建新的内部反例编号.
\end{FARemark}

\section{Riesz表示定理}\label{sec:i04-riesz-representation}
\index{Riesz representation theorem@Riesz表示定理}

本节只使用连续线性泛函 \(\displaystyle f:H\to\mathbb K\) 的拓扑定义，不建立 I-05 的一般有界线性算子或对偶空间理论.

\begin{FATheorem}[A][Hilbert空间Riesz表示定理]{fa:HIL-A02}
设 \(\displaystyle H\) 为 Hilbert 空间，\(\displaystyle f:H\to\mathbb K\) 为连续线性泛函. 则存在唯一 \(\displaystyle y_f\in H\)，使 \(\displaystyle \forall\,x\in H,\; f(x)=\langle x,y_f\rangle\).
并且 \(\displaystyle \sup_{\|x\|\leqslant1}|f(x)|=\|y_f\|\).
\end{FATheorem}

\begin{FAProof}
先设 \(\displaystyle f=0\). 取 \(\displaystyle y_f=0\)，则表示式成立；若另有 \(\displaystyle y\in H\) 满足 \(\displaystyle \langle x,y\rangle=0\) 对全部 \(\displaystyle x\in H\) 成立，取 \(\displaystyle x=y\) 得 \(\displaystyle \|y\|^2=0\)，故 \(\displaystyle y=0\). 此时上确界两端均为零.

现设 \(\displaystyle f\neq0\)，令 \(\displaystyle M=\ker f\). 若 \(\displaystyle x_n\in M\) 且 \(\displaystyle x_n\to x\)，则连续性给出 \(\displaystyle f(x_n)\to f(x)\)，而 \(\displaystyle f(x_n)=0\)，故 \(\displaystyle f(x)=0\)；因此 \(\displaystyle M\) 闭. 因 \(\displaystyle f\neq0\)，存在 \(\displaystyle z\in H\) 使 \(\displaystyle f(z)\neq0\)，所以 \(\displaystyle z\notin M\) 且 \(\displaystyle M\neq H\).

由\autoref{i04:closed-subspace-projection}，存在唯一分解 \(\displaystyle z=m_0+u, \; m_0\in M, \; u\in M^\perp\).
若 \(\displaystyle u=0\)，则 \(\displaystyle z=m_0\in M\)，矛盾，故 \(\displaystyle u\neq0\). 又 \(\displaystyle f(u)=f(z)-f(m_0)=f(z)\neq0\).

任取 \(\displaystyle x\in H\)，定义 \(\displaystyle m=x-\frac{f(x)}{f(u)}u\).
由线性性，\(\displaystyle f(m)=f(x)-f(x)=0\)，故 \(\displaystyle m\in M\). 因 \(\displaystyle u\in M^\perp\)，
\[
0=\langle m,u\rangle
=\langle x,u\rangle-\frac{f(x)}{f(u)}\|u\|^2,
\]
于是 \(\displaystyle \langle x,u\rangle =\frac{f(x)}{f(u)}\|u\|^2\).
在第一变量线性、第二变量共轭线性约定下，定义 \(\displaystyle y_f =\frac{\overline{f(u)}}{\|u\|^2}u\).
则对任意 \(\displaystyle x\in H\)，
\begin{align*}
\displaystyle \langle x,y_f\rangle
&\displaystyle =\overline{\frac{\overline{f(u)}}{\|u\|^2}}\langle x,u\rangle\\[12pt]
&\displaystyle =\frac{f(u)}{\|u\|^2}\cdot\frac{f(x)}{f(u)}\|u\|^2=f(x).
\end{align*}
故表示元存在.

若 \(\displaystyle y_1,y_2\in H\) 均满足 \(\displaystyle \langle x,y_1\rangle=\langle x,y_2\rangle\) 对全部 \(\displaystyle x\in H\) 成立，则 \(\displaystyle \langle x,y_1-y_2\rangle=0\). 取 \(\displaystyle x=y_1-y_2\)，得到 \(\displaystyle \|y_1-y_2\|^2=0\)，故 \(\displaystyle y_1=y_2\).

最后证明定量恒等式. Cauchy--Schwarz 给出 \(\displaystyle |f(x)| =|\langle x,y_f\rangle| \leqslant\|x\|\|y_f\|\)，
故 \(\displaystyle \sup_{\|x\|\leqslant1}|f(x)|\leqslant\|y_f\|\). 若 \(\displaystyle y_f=0\)，两边均为零. 若 \(\displaystyle y_f\neq0\)，取 \(\displaystyle x=\frac{y_f}{\|y_f\|}\)，则 \(\displaystyle \|x\|=1\) 且
\[
|f(x)|
=\left|\left\langle\frac{y_f}{\|y_f\|},y_f\right\rangle\right|
=\|y_f\|.
\]
故反向不等式成立，上确界恒等式得证.\hfill\(\square\)

\end{FAProof}

\section{正交基与Parseval恒等式}\label{sec:i04-orthonormal-bases}
\index{orthonormal basis@标准正交基}\index{Bessel inequality@Bessel不等式}\index{Parseval identity@Parseval恒等式}\index{Fourier coefficient@Fourier系数}

\subsection{任意指标集上的标准正交族}

\begin{FADefinition}[B][Hilbert标准正交基与任意指标求和]{i04:orthonormal-basis}
设 \(\displaystyle (e_j)_{j\in J}\) 为 \(\displaystyle H\) 中标准正交族. 若 \(\displaystyle \overline{\operatorname{span}}\{e_j:j\in J\}=H\)，
则称其为 Hilbert 意义的标准正交基. 这一定义不要求每个向量是有限线性组合，故一般不等同于 Hamel 基.

对非负族 \(\displaystyle (a_j)_{j\in J}\)，定义 \(\displaystyle \sum_{j\in J}a_j :=\sup_{F\Subset J}\sum_{j\in F}a_j\)，
其中 \(\displaystyle F\Subset J\) 表示 \(\displaystyle F\) 是 \(\displaystyle J\) 的有限子集. 对 \(\displaystyle x\in H\) 与有限 \(\displaystyle F\Subset J\)，定义 Fourier 有限部分和 \(\displaystyle S_Fx:=\sum_{j\in F}\langle x,e_j\rangle e_j\).
全体有限子集按包含关系构成有向集，\(\displaystyle (S_Fx)_{F\Subset J}\) 按此方向理解为网.
\end{FADefinition}

\begin{FATheorem}[B][Bessel不等式与Parseval]{fa:HIL-B01}
设 \(\displaystyle (e_j)_{j\in J}\) 为标准正交族，令 \(\displaystyle M:=\overline{\operatorname{span}}\{e_j:j\in J\}\).
则对每个 \(\displaystyle x\in H\)：
\begin{enumerate}[label=\textnormal{(\roman*)}]
\item Bessel 不等式成立： \(\displaystyle \sum_{j\in J}|\langle x,e_j\rangle|^2\leqslant\|x\|^2\);
\item Fourier 有限部分和网满足 \(\displaystyle S_Fx\to\mathcal P_Mx\)；
\item 有精确差额公式 \(\displaystyle \|x\|^2 =\|x-\mathcal P_Mx\|^2 +\sum_{j\in J}|\langle x,e_j\rangle|^2\).
特别地，若 \(\displaystyle (e_j)_{j\in J}\) 为标准正交基，则 \(\displaystyle S_Fx\to x, \; \|x\|^2=\sum_{j\in J}|\langle x,e_j\rangle|^2\).
\end{enumerate}
\end{FATheorem}

\begin{FAProof}
固定 \(\displaystyle x\in H\). 对任意有限 \(\displaystyle F\Subset J\) 与任意 \(\displaystyle k\in F\)，标准正交性给出
\[
\langle x-S_Fx,e_k\rangle
=\langle x,e_k\rangle
-\sum_{j\in F}\langle x,e_j\rangle\langle e_j,e_k\rangle
=0.
\]
故 \(\displaystyle x-S_Fx\perp S_Fx\). 由有限正交和公式， \(\displaystyle \|S_Fx\|^2 =\sum_{j\in F}|\langle x,e_j\rangle|^2\),
再由勾股恒等式， \(\displaystyle \|x\|^2 =\|x-S_Fx\|^2+\|S_Fx\|^2 \geqslant\sum_{j\in F}|\langle x,e_j\rangle|^2\).
对全部有限 \(\displaystyle F\Subset J\) 取上确界，即得 Bessel 不等式.

记 \(\displaystyle A=\sum_{j\in J}|\langle x,e_j\rangle|^2\)，则 \(\displaystyle A\leqslant\|x\|^2<\infty\). 给定 \(\displaystyle \varepsilon>0\)，由上确界定义可取有限 \(\displaystyle F_0\Subset J\) 使 \(\displaystyle A-\sum_{j\in F_0}|\langle x,e_j\rangle|^2<\varepsilon^2\).
若有限 \(\displaystyle F,G\Subset J\) 均包含 \(\displaystyle F_0\)，则利用标准正交性，
\begin{align*}
\|S_Fx-S_Gx\|^2
&=\sum_{j\in F\mathbin{\triangle}G}|\langle x,e_j\rangle|^2\\
&\leqslant\sum_{j\in(F\cup G)\setminus F_0}|\langle x,e_j\rangle|^2\\
&=\sum_{j\in F\cup G}|\langle x,e_j\rangle|^2
-\sum_{j\in F_0}|\langle x,e_j\rangle|^2\\
&\leqslant A-\sum_{j\in F_0}|\langle x,e_j\rangle|^2
<\varepsilon^2.
\end{align*}
故 \(\displaystyle (S_Fx)_{F\Subset J}\) 为 Cauchy 网. 现只用 Hilbert 空间的序列完备性证明该网收敛. 对每个 \(\displaystyle n\in\mathbb N\)，由上确界定义取有限 \(\displaystyle E_n\Subset J\) 使 \(\displaystyle A-\sum_{j\in E_n}|\langle x,e_j\rangle|^2<2^{-2n-2}\).
令 \(\displaystyle F_n=E_1\cup\dots\cup E_n\)，则 \(\displaystyle F_n\) 有限且递增，并且
\[
A-\sum_{j\in F_n}|\langle x,e_j\rangle|^2
\leqslant A-\sum_{j\in E_n}|\langle x,e_j\rangle|^2
<2^{-2n-2}.
\]
若 \(\displaystyle m\geqslant n\)，则 \(\displaystyle F_m\supseteq F_n\)，故 \(\displaystyle \|S_{F_m}x-S_{F_n}x\|^2 =\sum_{j\in F_m\setminus F_n}|\langle x,e_j\rangle|^2 <2^{-2n-2}\).
于是 \(\displaystyle (S_{F_n}x)\) 为 Cauchy 列. Hilbert 完备性给出 \(\displaystyle s\in H\) 使 \(\displaystyle S_{F_n}x\to s\). 给定 \(\displaystyle \varepsilon>0\)，取 \(\displaystyle n\) 足够大，使 \(\displaystyle 2^{-n-1}<\frac{\varepsilon}{2}\) 且 \(\displaystyle \|S_{F_n}x-s\|<\frac{\varepsilon}{2}\). 对任意有限 \(\displaystyle F\supseteq F_n\)， \(\displaystyle \|S_Fx-S_{F_n}x\|^2 \leqslant A-\sum_{j\in F_n}|\langle x,e_j\rangle|^2 <2^{-2n-2}\),
故 \(\displaystyle \|S_Fx-s\|<\varepsilon\). 因而整个有限子集网满足 \(\displaystyle S_Fx\to s\). 每个 \(\displaystyle S_Fx\in\operatorname{span}\{e_j:j\in J\}\)，而 \(\displaystyle M\) 闭，故 \(\displaystyle s\in M\).

固定 \(\displaystyle k\in J\). 当 \(\displaystyle F\supseteq\{k\}\) 时，已有 \(\displaystyle \langle x-S_Fx,e_k\rangle=0\). 由 Cauchy--Schwarz，内积对第一变量连续，因此令网取极限得 \(\displaystyle \langle x-s,e_k\rangle=0\). 线性性给出 \(\displaystyle x-s\) 与 \(\displaystyle \operatorname{span}\{e_j:j\in J\}\) 正交；再由同一连续性，它与其闭包 \(\displaystyle M\) 正交. 因而 \(\displaystyle s\in M\) 且 \(\displaystyle x-s\in M^\perp\). 由\autoref{i04:closed-subspace-projection}，\(\displaystyle s=\mathcal P_Mx\).

又 \(\displaystyle S_Fx\to\mathcal P_Mx\) 蕴含 \(\displaystyle \|S_Fx\|^2\to\|\mathcal P_Mx\|^2\). 另一方面，有限和网
\(\displaystyle \sum_{j\in F}|\langle x,e_j\rangle|^2\)
按定义递增到其上确界 \(\displaystyle A\). 因 \(\displaystyle \|S_Fx\|^2\) 正是该有限和，得到 \(\displaystyle \|\mathcal P_Mx\|^2 =\sum_{j\in J}|\langle x,e_j\rangle|^2\).
正交分解 \(\displaystyle x=\mathcal P_Mx+(x-\mathcal P_Mx)\) 再给出 \(\displaystyle \|x\|^2 =\|x-\mathcal P_Mx\|^2 +\sum_{j\in J}|\langle x,e_j\rangle|^2\).
若标准正交族为基，则 \(\displaystyle M=H\)，所以 \(\displaystyle \mathcal P_Mx=x\). 因而 \(\displaystyle S_Fx\to x\)，且差额项为零，得到 Parseval 恒等式.\hfill\(\square\)

\end{FAProof}

\begin{FATheorem}[B][正交基完备性判据]{fa:HIL-B02}
设 \(\displaystyle E=(e_j)_{j\in J}\) 为 \(\displaystyle H\) 中标准正交族. 以下条件等价：
\begin{enumerate}[label=\textnormal{(\arabic*)}]
\item \(\displaystyle \overline{\operatorname{span}}E=H\)；
\item \(\displaystyle E^\perp=\{0\}\)，其中 \(\displaystyle E^\perp=\{x\in H:\forall\,j\in J,\ \langle x,e_j\rangle=0\}\)；
\item 对每个 \(\displaystyle x\in H\)，Fourier 有限部分和网满足 \(\displaystyle S_Fx\to x\)；
\item 对每个 \(\displaystyle x\in H\)，成立 Parseval 恒等式 \(\displaystyle \|x\|^2=\sum_{j\in J}|\langle x,e_j\rangle|^2\).
\end{enumerate}
\end{FATheorem}

\begin{FAProof}
令 \(\displaystyle M=\overline{\operatorname{span}}E\). 先证明 \(\displaystyle E^\perp=M^\perp\). 若 \(\displaystyle x\in M^\perp\)，则每个 \(\displaystyle e_j\in M\)，故 \(\displaystyle x\in E^\perp\). 反之，若 \(\displaystyle x\in E^\perp\)，则第一变量固定时第二变量共轭线性，故 \(\displaystyle \langle x,m\rangle=0\) 对全部 \(\displaystyle m\in\operatorname{span}E\) 成立. 对 \(\displaystyle m_n\in\operatorname{span}E\) 且 \(\displaystyle m_n\to m\in M\)，Cauchy--Schwarz 给出 \(\displaystyle |\langle x,m-m_n\rangle|\leqslant\|x\|\|m-m_n\|\to0\)，故 \(\displaystyle \langle x,m\rangle=0\). 因而 \(\displaystyle x\in M^\perp\).

若（1）成立，则 \(\displaystyle M=H\)，故 \(\displaystyle M^\perp=\{0\}\)，于是（2）成立. 若（2）成立，则 \(\displaystyle M^\perp=\{0\}\). 由于 \(\displaystyle M\) 闭，\autoref{i04:closed-subspace-projection}给出 \(\displaystyle H=M\oplus M^\perp=M\)，故（1）成立. 因而（1）与（2）等价.

若（1）成立，则\autoref{fa:HIL-B01}中 \(\displaystyle S_Fx\to\mathcal P_Mx\) 化为 \(\displaystyle S_Fx\to x\)，故（3）成立. 若（3）成立，则每个 \(\displaystyle S_Fx\in M\)，且 \(\displaystyle M\) 闭，所以 \(\displaystyle x=\lim_FS_Fx\in M\) 对每个 \(\displaystyle x\in H\) 成立，故 \(\displaystyle M=H\)，即（1）成立.

若（3）成立，则对任意 \(\displaystyle x\in H\)，范数连续性与有限正交和公式给出
\[
\|x\|^2
=\lim_F\|S_Fx\|^2
=\sup_{F\Subset J}\sum_{j\in F}|\langle x,e_j\rangle|^2
=\sum_{j\in J}|\langle x,e_j\rangle|^2,
\]
故（4）成立. 若（4）成立且 \(\displaystyle x\in E^\perp\)，则每个 Fourier 系数均为零，于是 \(\displaystyle \|x\|^2=0\)，故 \(\displaystyle x=0\). 因而 \(\displaystyle E^\perp=\{0\}\)，即（2）成立. 四个条件因此等价.\hfill\(\square\)

\end{FAProof}

\subsection{\texorpdfstring{\(\displaystyle\ell^2\)}{ℓ^2}标准基与正交展开}

\begin{FAProposition}[B][\(\displaystyle\ell^2\)的标准单位向量基]{i04:ell2-standard-basis}
令 \(\displaystyle e_n=(0,\dots,0,1,0,\dots)\in\ell^2\). 则 \(\displaystyle (e_n)_{n\in\mathbb N}\) 是 \(\displaystyle \ell^2\) 的标准正交基. 对 \(\displaystyle x=(x_n)\in\ell^2\)，
\[
x=\lim_{N\to\infty}\sum_{n=1}^{N}\langle x,e_n\rangle e_n,
\quad
\|x\|_2^2=\sum_{n=1}^{\infty}|\langle x,e_n\rangle|^2=\sum_{n=1}^{\infty}|x_n|^2.
\]
\end{FAProposition}

\begin{FAProof}
由\autoref{i04:ell2-hilbert}，\(\displaystyle \langle e_n,e_m\rangle_{\ell^2}=0\) 当 \(\displaystyle n\neq m\)，且 \(\displaystyle \|e_n\|_2=1\)，故该族标准正交. 对 \(\displaystyle x=(x_n)\in\ell^2\)，有 \(\displaystyle \langle x,e_n\rangle=x_n\). 定义 \(\displaystyle S_Nx=(x_1,\dots,x_N,0,\dots) =\sum_{n=1}^{N}x_ne_n\).
则 \(\displaystyle \|x-S_Nx\|_2^2 =\sum_{n>N}|x_n|^2\to0\)，
因为 \(\displaystyle \sum_{n=1}^{\infty}|x_n|^2\) 收敛. 因而有限线性组合在 \(\displaystyle \ell^2\) 中稠密，\(\displaystyle (e_n)\) 是标准正交基. 所列展开与 Parseval 公式分别由该收敛和\autoref{fa:HIL-B01}得到.\hfill\(\square\)

\end{FAProof}

\begin{FAApplication}[抽象正交展开]
对任意 Hilbert 标准正交基 \(\displaystyle (e_j)_{j\in J}\)，每个 \(\displaystyle x\in H\) 都由其 Fourier 系数族 \(\displaystyle (\langle x,e_j\rangle)_{j\in J}\) 通过有限部分和网恢复：\(\displaystyle S_Fx\to x\). Parseval 同时给出其范数平方等于 Fourier 系数模平方族的任意指标和. 在 \(\displaystyle \ell^2\) 中，这一抽象展开退化为坐标截断收敛.
\end{FAApplication}

\subsection{本章习题}\label{sec:i04-exercises}

\begin{FAExercise}[内积公理与复共轭位置]{ex:i04-01}
在 \(\displaystyle \mathbb C^2\) 上定义 \(\displaystyle \langle x,y\rangle=x_1\overline{y_1}+2x_2\overline{y_2}\). 逐项验证第一变量线性、第二变量共轭线性、共轭对称性与正定性，并写出诱导范数.
\end{FAExercise}

\begin{FAExercise}[Cauchy--Schwarz中的最优参数]{ex:i04-02}
对 \(\displaystyle y\neq0\)，从 \(\displaystyle \|x-\lambda y\|^2\geqslant0\) 出发，将右端写成关于 \(\displaystyle \lambda\) 的完整表达式，并证明 \(\displaystyle \lambda=\frac{\langle x,y\rangle}{\|y\|^2}\) 恰使余项最小. 不得使用三角不等式.
\end{FAExercise}

\begin{FAExercise}[Cauchy--Schwarz等号条件]{ex:i04-03}
分别处理 \(\displaystyle x=0\)、\(\displaystyle y=0\) 与 \(\displaystyle x,y\neq0\)，完整证明 \(\displaystyle |\langle x,y\rangle|=\|x\|\|y\|\) 当且仅当 \(\displaystyle x,y\) 线性相关.
\end{FAExercise}

\begin{FAExercise}[复极化恒等式符号核验]{ex:i04-04}
在第一变量线性约定下直接展开 \(\displaystyle \|x\pm\mathrm{i}y\|^2\)，由此重新推导\autoref{fa:HIL-C01}中的复极化恒等式，并解释若误用第二变量线性公式会在哪一项改变符号.
\end{FAExercise}

\begin{FAExercise}[正交补的闭性]{ex:i04-05}
对任意集合 \(\displaystyle A\subseteq H\)，定义 \(\displaystyle A^\perp=\{x:\forall\,a\in A,\ \langle x,a\rangle=0\}\). 证明 \(\displaystyle A^\perp\) 是闭线性子空间，并说明证明只需 Cauchy--Schwarz 与内积线性.
\end{FAExercise}

\begin{FAExercise}[有限正交和恒等式]{ex:i04-06}
设 \(\displaystyle x_1,\dots,x_n\) 两两正交. 不使用归纳法，直接展开 \(\displaystyle \langle\sum_kx_k,\sum_jx_j\rangle\)，证明 \(\displaystyle \|\sum_kx_k\|^2=\sum_k\|x_k\|^2\).
\end{FAExercise}

\begin{FAExercise}[闭凸投影的存在性估计]{ex:i04-07}
重做\autoref{fa:HIL-A01}的存在性部分. 从 \(\displaystyle \|x-y_n\|^2\to d^2\) 与凸性出发，给定 \(\displaystyle \varepsilon>0\) 显式选择 \(\displaystyle N\)，使 \(\displaystyle n,m\geqslant N\Rightarrow\|y_n-y_m\|<\varepsilon\).
\end{FAExercise}

\begin{FAExercise}[投影唯一性与严格凸性]{ex:i04-08}
只使用平行四边形恒等式与中点凸性证明闭凸集最佳逼近元唯一，不得把 Hilbert 范数严格凸作为已知定理. 随后从同一计算推出内积范数单位球的严格凸性.
\end{FAExercise}

\begin{FAExercise}[变分不等式方向]{ex:i04-09}
从 \(\displaystyle \|x-y\|^2\leqslant\|x-y-t(z-y)\|^2\)、\(\displaystyle t\in(0,1]\) 完整展开并除以 \(\displaystyle t\)，证明极限给出的是 \(\displaystyle \operatorname{Re}\langle x-y,z-y\rangle\leqslant0\) 而非反向不等式.
\end{FAExercise}

\begin{FAExercise}[复闭子空间的正交条件]{ex:i04-10}
在复 Hilbert 空间中，设 \(\displaystyle \operatorname{Re}\langle v,w\rangle=0\) 对全部 \(\displaystyle w\in M\) 成立，其中 \(\displaystyle M\) 是复线性子空间. 分别代入 \(\displaystyle w\) 与 \(\displaystyle \mathrm{i}w\)，证明 \(\displaystyle v\in M^\perp\).
\end{FAExercise}

\begin{FAExercise}[投影的线性与幂等性]{ex:i04-11}
设 \(\displaystyle M\leqslant H\) 闭. 在不引入一般算子范数理论的前提下，由正交分解唯一性证明 \(\displaystyle \mathcal P_M\) 线性，并证明 \(\displaystyle \mathcal P_M(\mathcal P_Mx)=\mathcal P_Mx\) 与 \(\displaystyle \mathcal P_Mx=x\iff x\in M\).
\end{FAExercise}

\begin{FAExercise}[双正交补的非闭情形]{ex:i04-12}
设 \(\displaystyle M\leqslant H\) 不闭. 证明 \(\displaystyle M\subsetneq M^{\perp\perp}\)，并给出等号 \(\displaystyle M^{\perp\perp}=\overline M\) 中两次使用闭性的准确位置.
\end{FAExercise}

\begin{FAExercise}[有限正交最小二乘]{ex:i04-13}
设 \(\displaystyle e_1,\dots,e_n\) 标准正交. 对任意标量 \(\displaystyle c_1,\dots,c_n\)，直接证明
\[
\left\|x-\sum_{k=1}^{n}c_ke_k\right\|^2
=\left\|x-\sum_{k=1}^{n}\langle x,e_k\rangle e_k\right\|^2
+\sum_{k=1}^{n}|c_k-\langle x,e_k\rangle|^2,
\]
并由此重新得到唯一最小二乘系数.
\end{FAExercise}

\begin{FAExercise}[Riesz证明中的核闭性]{ex:i04-14}
设 \(\displaystyle f:H\to\mathbb K\) 连续线性且 \(\displaystyle f\neq0\). 只用连续性的序列定义证明 \(\displaystyle \ker f\) 闭，并证明任取 \(\displaystyle z\notin\ker f\) 后其正交分量 \(\displaystyle u\) 必满足 \(\displaystyle u\neq0\) 与 \(\displaystyle f(u)\neq0\).
\end{FAExercise}

\begin{FAExercise}[Riesz复系数专项]{ex:i04-15}
沿用\autoref{fa:HIL-A02}证明中的 \(\displaystyle u\). 分别计算候选 \(\displaystyle y_1=\frac{f(u)u}{\|u\|^2}\) 与 \(\displaystyle y_2=\frac{\overline{f(u)}u}{\|u\|^2}\) 所产生的 \(\displaystyle \langle x,y_j\rangle\)，并判断哪个候选满足 \(\displaystyle f(x)=\langle x,y_j\rangle\).
\end{FAExercise}

\begin{FAExercise}[Riesz的上确界恒等式]{ex:i04-16}
已知 \(\displaystyle f(x)=\langle x,y\rangle\). 不使用 I-05 的算子范数语言，分别处理 \(\displaystyle y=0\) 与 \(\displaystyle y\neq0\)，证明 \(\displaystyle \sup_{\|x\|\leqslant1}|f(x)|=\|y\|\).
\end{FAExercise}

\begin{FAExercise}[任意指标Bessel求和]{ex:i04-17}
设 \(\displaystyle (e_j)_{j\in J}\) 标准正交. 从每个有限 \(\displaystyle F\Subset J\) 的勾股恒等式出发，证明 \(\displaystyle \sum_{j\in J}|\langle x,e_j\rangle|^2\leqslant\|x\|^2\). 不得预设 \(\displaystyle J\) 可数.
\end{FAExercise}

\begin{FAExercise}[Fourier有限子集网的Cauchy性]{ex:i04-18}
设 \(\displaystyle A=\sup_{F\Subset J}\sum_{j\in F}|\langle x,e_j\rangle|^2\). 给定 \(\displaystyle \varepsilon>0\)，从上确界定义选取有限 \(\displaystyle F_0\)，并完整证明 \(\displaystyle F,G\supseteq F_0\Rightarrow\|S_Fx-S_Gx\|<\varepsilon\).
\end{FAExercise}

\begin{FAExercise}[Parseval与闭线性张成]{ex:i04-19}
对标准正交族 \(\displaystyle E\) 与 \(\displaystyle M=\overline{\operatorname{span}}E\)，证明对固定 \(\displaystyle x\in H\)，Parseval 等号成立当且仅当 \(\displaystyle x\in M\). 必须使用\autoref{fa:HIL-B01}的差额公式.
\end{FAExercise}

\begin{FAExercise}[正交基完备性判据]{ex:i04-20}
不引用\autoref{fa:HIL-B02}的结论，独立证明 \(\displaystyle E^\perp=\{0\}\Rightarrow\overline{\operatorname{span}}E=H\). 要求显式使用闭子空间正交分解，而不是只写“双正交补即得”.
\end{FAExercise}

\begin{FAExercise}[\(\displaystyle\ell^2\)截断与Parseval]{ex:i04-21}
对 \(\displaystyle x=(x_n)\in\ell^2\)，证明 \(\displaystyle S_Nx=(x_1,\dots,x_N,0,\dots)\) 满足 \(\displaystyle \|x-S_Nx\|_2^2=\sum_{n>N}|x_n|^2\to0\)，并直接核验 Fourier 系数 \(\displaystyle \langle x,e_n\rangle=x_n\).
\end{FAExercise}

\begin{FAExercise}[\(\displaystyle L^2\)代表元与最小二乘]{ex:i04-22}
设 \(\displaystyle g_1,\dots,g_n\in L^2(\Omega)\) 标准正交. 先证明 \(\displaystyle \langle f,g_k\rangle_{L^2}\) 与代表元选择无关，再证明 \(\displaystyle \sum_{k=1}^{n}\langle f,g_k\rangle_{L^2}g_k\) 是 \(\displaystyle f\) 在 \(\displaystyle \operatorname{span}\{g_1,\dots,g_n\}\) 中的唯一最佳逼近.
\end{FAExercise}
